\section{Overview}

\subsection{Problem Statement}

On university campuses, students still organize study groups mostly by hand, through disconnected chat groups on Facebook, Zalo, or Messenger. This approach has three concrete shortcomings from a software-engineering standpoint:

\begin{itemize}
    \item There is no way to filter by course, expectation level, or physical proximity within a campus, so a student has to manually sift through dozens of unrelated chat groups.
    \item There is no reliable way to tell whether a group still has an open seat, is already full, or contains a member the student would rather not study with.
    \item There is no activity history, no identity guarantee that the organizer actually belongs to the claimed school/department, and no automated reminder mechanism.
\end{itemize}

\textbf{Study Group Finder} (in-app display name: \emph{StudyCohort}) is an Android application built to replace this ad-hoc workflow with a structured system: students who belong to the same \emph{community} (school/department) can create, discover, and join study sessions filtered by course, time, location, and commitment level.

\subsection{Objectives}

\begin{itemize}
    \item Build a complete Android application following the MVVM pattern, using Firebase (Authentication, Firestore, built-in offline persistence) as the primary data source and Room as the local cache layer.
    \item Support the full lifecycle of a study session: creation, member invitation, join-request approval, editing, cancellation, and completion.
    \item Provide session discovery through filtering, search, and sorting (by time, name, and geographic distance via the Haversine formula combined with device GPS).
    \item Satisfy all three mandatory technical requirements of the course (Section~\ref{sec:coursereq}).
    \item Deliver a consistent user experience with all four standard UI states (loading, empty, error, offline) on every screen that fetches data.
\end{itemize}

\subsection{Scope and Target Users}

The application follows a multi-community model: after registering, each user selects (or is routed to) exactly one \emph{community} — representing a specific school or department — and only ever sees sessions, members, and campus locations that belong to that community. A community can be \emph{verified} (requires an email address on the school's domain, plus a verified email) or \emph{unverified} (open to anyone).

The target audience is university students who need to find or organize course-based study groups. The technical scope is limited to Android (no iOS or web client), with no dedicated backend server — all business logic runs on the client and is constrained by Firestore Security Rules — and no server-triggered push notifications (FCM); reminders are handled entirely on-device via WorkManager.

\subsection{Mandatory Course Requirements}
\label{sec:coursereq}

The project must independently demonstrate three core technical capabilities required by the course, separate from any business feature. Table~\ref{tab:coursereq} summarizes how each requirement is satisfied and where it is implemented.

\begin{table}[h]
\centering
\caption{The three mandatory technical requirements and how they are satisfied}
\label{tab:coursereq}
\begin{tabular}{p{3.2cm}p{6cm}p{4cm}}
\toprule
\textbf{Requirement} & \textbf{How it is satisfied} & \textbf{Source location} \\
\midrule
Persistent local storage & Room (SQLite) acts as a read-only offline cache for every core screen; data remains visible with the network off & \texttt{data/local/} (4 entities, 4 DAOs) \\
\addlinespace
External API call & Retrofit calls \texttt{https://firestore.\allowbreak googleapis.com/v1/\ldots} (the raw Firestore REST API) directly to load the initial community list, bypassing the Firebase SDK & \texttt{data/remote/rest/PublicCommunityApi.kt} \\
\addlinespace
Device hardware capability & Camera (profile photo capture via a system intent) and GPS (\texttt{FusedLocationProviderClient} for distance-based sorting) & \texttt{ui/profile/ProfileFragment.kt}, \texttt{ui/home/HomeFragment.kt} \\
\bottomrule
\end{tabular}
\end{table}

The second requirement (external API) is deliberately implemented as a manual REST call rather than the Firestore SDK, even though both reach the same underlying database — this is precisely the course's \emph{external API} requirement, and silently swapping it for the SDK would no longer satisfy it. The response uses Firestore's REST \emph{typed wire format} (every field wrapped in an object such as \texttt{\{"stringValue": ...\}}), which requires a dedicated DTO layer (\texttt{dto/FirestoreRestDto.kt}) before being mapped to the domain model.

\subsection{Team Responsibilities}

Work was split along two axes: account/infrastructure ownership, and session-lifecycle ownership. Table~\ref{tab:assignment} summarizes the primary assignment; the remaining screens (Home, Community, Inbox) were implemented jointly, since they depend on both areas.

\begin{table}[h]
\centering
\caption{Detailed work assignment by module}
\label{tab:assignment}
\begin{tabular}{p{2.3cm}p{9.5cm}p{1.6cm}}
\toprule
\textbf{Member} & \textbf{Module / screen ownership} & \textbf{\%} \\
\midrule
Khoi &
Firebase Console setup (Authentication, Firestore, Storage) and Security Rules authoring; \newline
Shared navigation graph (\texttt{nav\_graph.xml}) and app header component; \newline
Splash screen and three-way routing logic; \newline
Seed-data script (\texttt{DataSeeder}); \newline
Sign up / sign in / forgot password / sign out; \newline
Profile screen (edit fields, activity graph, photo upload, block user, dual view modes)
& \\
\addlinespace
Vi &
Create / Edit / Manage session (including the staged-changes pattern); \newline
Session Detail (action-state machine, study material attachments); \newline
My Sessions (list view and self-built calendar view); \newline
Session History (including PDF export)
& \\
\addlinespace
Shared &
Home (filters, search, sort, Haversine distance); \newline
Community Selection (REST API, verified gate); \newline
Inbox (three notification types, WorkManager reminders); \newline
The four UI states applied across the whole application
& \\
\bottomrule
\end{tabular}
\end{table}

\emph{The \% column is left for the team to fill in based on an internal contribution assessment before submission.}
